\chapter{The Technics message set}\label{ch:messageset}

\section{Message families}
After \bytes{F0 50}, the third byte selects a message family. The instrument accepts
fourteen values:

\begin{center}
\bytes{21 22 23 24 25 27 28 29 2A 2B 2C 2D 7E 7F}
\end{center}

\bytes{26} is not accepted. The SX-WSA1R's own product identifier is \bytes{2D}; its
outgoing bulk-dump data messages all begin \bytes{F0 50 2D}.

\section{The model triple}
The families \bytes{21}, \bytes{22}, \bytes{2B}, \bytes{2C} and \bytes{2D} continue with
three fixed bytes:

\begin{center}
\bytes{04} \quad \bytes{\textit{nn}} \quad \bytes{11}
\end{center}

where \bytes{\textit{nn}} is \bytes{00} or \bytes{01}. A message whose triple the
instrument does not recognise is rejected with \textsc{error 40}.

\begin{caution}{The middle byte names the model, and only on transmission}
The keyboard transmits \bytes{04 00 11}; the rack transmits \bytes{04 01 11}. The rack
substitutes the middle byte as the message leaves, in every \bytes{21}, \bytes{22},
\bytes{2C} and \bytes{2D} message it sends, and recomputes the checksum of those that
carry one, so the substitution is invisible in the checksum. A bulk dump captured from a
real rack carries \bytes{04 01 11} throughout.

Reception is model-blind. Both values are accepted by both models, on every family that
carries the triple, and reach the same handler; a message is never refused, ignored or
treated differently because of the middle byte. A receiver must therefore accept both,
and a file dumped from one model loads into the other.
\end{caution}

\section{Short messages}
These messages are complete as shown. They form the control vocabulary of a bulk-dump
session.

\begin{center}
\begin{tabular}{ll}
\toprule
message & meaning \\
\midrule
\bytes{F0 50 23 7E F7} & acknowledge \\
\bytes{F0 50 24 7E F7} & not acknowledged \\
\bytes{F0 50 27 7E F7} & end of category \\
\bytes{F0 50 28 7E F7} & end of dump \\
\bytes{F0 50 29 7E F7} & abort \\
\bytes{F0 50 2A 7E F7} & destination full \\
\bytes{F0 50 21 04 \textit{nn} 11 F7} & enquiry \\
\bytes{F0 50 22 04 \textit{nn} 11 F7} & enquiry reply \\
\bottomrule
\end{tabular}
\end{center}

\section{Tempo}
The family \bytes{25} carries a tempo setting, in the range 40 to 300 beats per minute.
A value outside that range is treated as 120.

See chapter~\ref{ch:models} --- the tempo message is one of two things the rack does not
support.

\section{Parameter families}
The families \bytes{2B} and \bytes{2C} carry individual parameter traffic rather than bulk
transfer. Each addresses a large parameter space, and neither is enumerated in this
edition; see chapter~\ref{ch:limits}.
